% =============================================================================
% Dokumentace projektu IFJ/IAL 2026/2027 -- překladač IFJ26 -> IFJcode26
% Struktura a styl vycházejí z dokumentace z roku 2020.
%
% Překlad:   make            (latexmk -pdf; vyžaduje českou lokalizaci babel)
% Tabulky:   make tabulky    (gen_tabulky.py přegeneruje *.tex tabulky)
%
% Místa, která musí doplnit tým, jsou \todo{...} (červeně). Před odevzdáním
% je smažte a přepněte \drafttrue na \draftfalse.
% =============================================================================
\documentclass[11pt]{article}

\usepackage[T1]{fontenc}
\usepackage[utf8]{inputenc}
\usepackage[czech]{babel}
\usepackage{mathptmx}                       % písmo Times
\usepackage[a4paper, top=3cm, left=2cm, textwidth=17cm, textheight=24cm]{geometry}
\usepackage{parskip}
\usepackage{amsmath,amssymb}
\usepackage{graphicx}
\usepackage{xcolor}
\usepackage{booktabs,tabularx,longtable,array}
\usepackage{enumitem}
\usepackage{caption}
\usepackage{tikz}
\usetikzlibrary{arrows.meta,positioning}
\usepackage{pdflscape}
\usepackage[unicode,pdftitle={Překladač pro jazyk IFJ26},hidelinks]{hyperref}

\setlist{itemsep=1pt,topsep=3pt,parsep=0pt}
\captionsetup{font=small,labelfont=bf}

% --- značky pro nedodělaná místa --------------------------------------------
\newif\ifdraft
\drafttrue
\newcommand{\todo}[1]{\ifdraft{\color{red}\textbf{[TODO: #1]}}\fi}

% --- sazba -------------------------------------------------------------------
\newcommand{\nt}[1]{$\langle$\textit{#1}$\rangle$}   % neterminál
\newcommand{\file}[1]{\texttt{#1}}                   % soubor, identifikátor
\newcommand{\ic}[1]{\texttt{#1}}                     % instrukce IFJcode26
\newcommand{\dq}{\textquotedbl{}}                    % znak "
\newcommand{\cf}[1]{\textcolor{red!70!black}{\textbf{#1}}} % konflikt v LL tabulce

\begin{document}

% =============================================================================
% Titulní strana
% =============================================================================
\begin{titlepage}
  \centering
  \IfFileExists{img/FIT_barevne_CMYK_CZ.pdf}%
    {\includegraphics[width=\textwidth]{img/FIT_barevne_CMYK_CZ.pdf}\par}%
    {\todo{logo FIT: img/FIT\_barevne\_CMYK\_CZ.pdf}\par}
  \vspace{\stretch{0.382}}

  {\huge\bfseries Překladač pro jazyk IFJ26 \par}
  {\Large \textbf{Dokumentace k projektu} \par}
  {\Large Tým \todo{číslo týmu}, varianta vv-BVS \par}
  \vspace{1.5cm}

  \large
  \begin{tabular}{l l l l | r}
    \textbf{\todo{Jméno Příjmení}} & (\todo{xlogin00}) & \todo{os.\,č.} & -- \textbf{vedoucí} & \todo{25\,\%} \\
    \todo{Matěj Okruhlica} & (\todo{xokruhm00}) & \todo{os.\,č.} & & \todo{25\,\%} \\
    \todo{Alex Oľšavský} & (\todo{xolsava00}) & \todo{os.\,č.} & & \todo{25\,\%} \\
    \todo{Jaroslav Vančura} & (\todo{xvancuj00}) & \todo{os.\,č.} & & \todo{25\,\%} \\
    \todo{Klára Vantuchová} & (\todo{xvantuk00}) & \todo{os.\,č.} & & \todo{25\,\%} \\
  \end{tabular}

  \vspace{1.2cm}
  {\Large Implementovaná rozšíření: \par}
  {\large FUNEXP, BOOLTHEN, CONTROL, VALMEM, INTERPOLATEDS, EXTKW, PARAMLISTS \par}

  \vspace{0.8cm}
  {\large Git repozitář: \todo{odkaz} \par}

  \vspace{\stretch{0.618}}
  {\large \today\par}
\end{titlepage}

\tableofcontents
\newpage

% =============================================================================
\section{Práce v týmu}
% =============================================================================
\todo{stručně: jak jsme si práci naplánovali, jak se plnily termíny, kdy jsme spolupracovali (rozhraní mezi moduly, párové programování)}

\subsection{Rozdělení práce}

\begin{itemize}
  \item \todo{Jméno (login)}
    \begin{itemize}
      \item \todo{např.\ lexikální analýza, diagram konečného automatu, \ldots}
    \end{itemize}
  \item \todo{Jméno (login)}
    \begin{itemize}
      \item \todo{např.\ syntaktická analýza, precedenční analýza, LL gramatika a tabulka, \ldots}
    \end{itemize}
  \item \todo{Jméno (login)}
    \begin{itemize}
      \item \todo{např.\ tabulka symbolů, sémantické kontroly, \ldots}
    \end{itemize}
  \item \todo{Jméno (login)}
    \begin{itemize}
      \item \todo{např.\ generátor cílového kódu, testovací sada, dokumentace, \ldots}
    \end{itemize}
\end{itemize}

\todo{povinně zdůvodnit odchylky od rovnoměrného rozdělení bodů, pokud nastaly}

\subsection{Organizace práce}

\todo{verzovací systém, větve, pull requesty, kolik schválení se vyžaduje, kultura code review}

Testování probíhalo ve třech formách. \emph{End-to-end testy} přeloží zdrojový program IFJ26 překladačem, vygenerovaný kód provedou v interpretu \file{ic26int} a porovnají standardní výstup i návratový kód s očekávaným; testy jsou rozdělené do skupin podle částí jazyka (adresáře \file{tests/generator} a \file{tests/compiler}). Očekávané výstupy jsme ověřovali překladačem Scala (\file{ifj26.scala}). \emph{Statické kontroly} vygenerovaného kódu hlídají unikátnost návěští, existenci cílů skoků a to, že se \ic{DEFVAR} vypisují jen na začátku funkce. \todo{další formy testování}

Kód formátujeme nástrojem \todo{např.\ clang-format, konfigurace}; funkce jsou dokumentovány komentáři ve formátu Doxygen (\texttt{@brief}, \texttt{@param}).

% =============================================================================
\section{Návrh překladače}
% =============================================================================
Náš překladač čte zdrojový program ze standardního vstupu a vypisuje kód IFJcode26 na standardní výstup. Skládá se z lexikální analýzy, syntaktické a sémantické analýzy, tabulky symbolů a generátoru kódu (obrázek~\ref{fig:arch}). Syntaktická analýza je řízená rekurzivním sestupem a výrazy zpracovává precedenční syntaktickou analýzou. Překlad je syntaxí řízený: parser staví abstraktní syntaktický strom (AST) anotovaný typy výrazů a z něj generátor vytváří výstupní kód. Při chybě překladač skončí s příslušným návratovým kódem a na výstup nevypíše nic.

\begin{figure}[ht]
\centering
\begin{tikzpicture}[
    >=Stealth, node distance=11mm,
    blk/.style={draw, rounded corners=3pt, minimum width=4.2cm, minimum height=10mm,
                align=center, fill=gray!8, font=\small},
    io/.style={draw, minimum width=4.2cm, minimum height=8mm, align=center,
               font=\small, fill=white},
    dat/.style={font=\footnotesize\itshape, align=left},
    err/.style={draw, dashed, rounded corners=3pt, minimum width=3cm,
                minimum height=10mm, align=center, font=\small}
  ]
  \node[io]  (src) {Zdrojový kód IFJ26\\(\texttt{stdin})};
  \node[blk, below=of src] (lex) {Lexikální analýza\\\file{lexer.c}};
  \node[blk, below=of lex, minimum width=5.4cm] (par)
        {Syntaktická a sémantická analýza\\\file{parser.c}, \file{precedence.c}};
  \node[blk, right=34mm of par, minimum width=3.6cm] (sym)
        {Tabulka symbolů\\\file{symtable.c} (AVL)};
  \node[blk, below=of par] (gen) {Generátor kódu\\\file{generator.c}};
  \node[io, below=of gen] (out) {Kód IFJcode26\\(\texttt{stdout})};
  \node[err, left=16mm of par] (err) {Chyba:\\návratový kód\\\file{Error.h}};

  \draw[->] (src) -- (lex) node[midway, right, dat] {posloupnost znaků};
  \draw[->] (lex) -- (par) node[midway, right, dat] {tokeny (\texttt{Token})};
  \draw[<->] (par) -- (sym) node[midway, above, dat] {záznamy\\[-2pt]\texttt{symtable\_entry}};
  \draw[->] (par) -- (gen) node[midway, right, dat]
        {AST (\texttt{ASTNode}) s typy\\[-2pt]výrazů (\texttt{expr\_type})};
  \draw[->] (gen) -- (out) node[midway, right, dat] {textový kód};
  \draw[->, dashed] (lex.west) -| (err.north);
  \draw[->, dashed] (par.west) -- (err.east);
  \draw[->, dashed] (gen.west) -| (err.south);
\end{tikzpicture}
\caption{Blokové schéma překladače a datové toky mezi moduly.}
\label{fig:arch}
\end{figure}

\subsection{Lexikální analýza}
\label{sec:lexer}

Lexikální analyzátor (\file{lexer.c}) je implementovaný jako deterministický konečný automat, jehož diagram je v \hyperref[sec:fa]{příloze}~\ref{sec:fa}. Funkce \texttt{get\_next\_token} čte vstup po znacích a vrací vždy jeden token: jeho typ a podle typu lexém (identifikátory, literály) nebo počet mezer (odsazení). Lexém se při čtení skládá do dynamického řetězce (\file{DynBuff}) a do tokenu se předává jeho kopie. Že identifikátor je klíčové slovo, rozhoduje \texttt{keyword\_check} po načtení celého identifikátoru. Lexikální chyba se vrací jako token \texttt{T\_ERROR}, který parser převede na návratový kód~1.

Odchylky od standardního modelu konečného automatu:
\begin{itemize}
  \item \textbf{Odsazení a konce řádků.} Bloky jsou v IFJ26 vymezeny odsazením, proto lexer vrací počet mezer na začátku logického řádku jako token \texttt{T\_INDENTS} a konec řádku jako \texttt{T\_EOL}. Tabulátor na začátku řádku je lexikální chyba.
  \item \textbf{Vnořené blokové komentáře} automat bez čítače nezvládne, proto se hloubka vnoření počítá. Komentář obsahující konec řádku se chová jako konec řádku.
  \item \textbf{Řetězce.} U běžných řetězců se escape sekvence (\texttt{\textbackslash n}, \texttt{\textbackslash r}, \texttt{\textbackslash t}, \texttt{\textbackslash\textbackslash}, \texttt{\textbackslash\dq}) převádějí přímo v lexeru na odpovídající znak. Víceřádkové řetězce (\texttt{\dq\dq\dq}) se čtou doslovně.
  \item \textbf{Čísla.} Podtržítka mezi číslicemi zůstávají v lexému a odstraňují se až při převodu na hodnotu. Desetinný literál může končit příponou \texttt{d} nebo \texttt{D}.
  \item \textbf{Anotace \ic{@main}} je jeden token (\texttt{T\_FUNC\_MAIN}).
\end{itemize}

\subsection{Syntaktická analýza}

\subsubsection{Rekurzivní sestup}

Parser (\file{parser.c}) implementuje rekurzivní sestup podle LL gramatiky v příloze~\ref{sec:ll-gram}. Tokeny z lexeru se před zpracováním ukládají do bufferu (\texttt{raw\_token}), takže parser může vstup procházet opakovaně a nahlížet na následující token (\texttt{peek\_token}). Konec bloku se pozná podle odsazení: gramatika s ním pracuje jako s virtuálním terminálem \textit{DEDENT}.

Gramatika je LL(1) až na čtyři místa (tabulka v příloze~\ref{sec:ll-tab}): rozlišení přiřazení od výrazu začínajícího identifikátorem (rozhodne druhý token, \texttt{=}), příkaz \texttt{if} oproti výrazu \texttt{if} na začátku příkazu, a dvě místa týkající se větve \texttt{else} (nejbližší \texttt{if} a zarovnání \texttt{else} na stejné odsazení jako \texttt{if}). Každé řešíme kontextem, ne změnou gramatiky.

Zadání dovoluje volat funkci před její definicí a funkce se mohou volat vzájemně. Z toho důvodu parser před hlavním průchodem projde buffer tokenů a zapíše do globální tabulky symbolů hlavičky všech funkcí (\texttt{collect\_signatures}); těla funkcí přitom přeskočí. Tělo funkce se pak analyzuje při hlavním průchodu v pořadí zdrojového kódu (\texttt{analyze\_function}), a je-li při volání potřeba návratový typ funkce, jejíž tělo ještě nebylo zpracováno, analyzuje se její tělo přednostně (\texttt{resolve\_function}); stav analýzy těla si funkce nese v tabulce symbolů. \todo{ověřit popis proti precedence.c}

\subsubsection{Precedenční syntaktická analýza}
\label{sec:prec}

Výrazy zpracovává precedenční analýza (\file{precedence.c}) se zásobníkem operátorů a zásobníky typů a uzlů AST. Při redukci operátoru se z operandů vytvoří uzel AST a vypočítá se typ výsledku (včetně implicitního převodu Int na Double); nesouhlasí-li typy operandů, analýza skončí typovou chybou. Unární a binární \texttt{-} se rozliší podle toho, zda se na daném místě očekává operand, nebo operátor. Volání funkce uvnitř výrazu zpracovává rekurzivně \texttt{parse\_call} a z pohledu precedenční analýzy se chová jako operand. Priority a asociativita operátorů, včetně operátorů rozšíření BOOLTHEN, jsou v příloze~\ref{sec:prec-tab}.

\subsubsection{Sémantické kontroly}

Sémantické kontroly provádí parser za běhu analýzy. Kontroluje se existence a redefinice proměnných a funkcí, počet a typy argumentů při volání, kompatibilita typů ve výrazech, přiřazeních a inicializacích, zápis do nemodifikovatelné proměnné nebo parametru a pravidla pro příkaz \texttt{return}. Návratový typ funkce bez explicitního typu se odvodí z jejího těla. Existenci právě jedné funkce s anotací \ic{@main}, bez parametrů a s návratovým typem \texttt{Unit}, kontroluje parser při zpracování definic funkcí. Typ každého výrazu se ukládá do uzlu AST (\texttt{expr\_type}), takže generátor už žádné sémantické kontroly neprovádí.

\subsection{Vnitřní reprezentace kódu}

Parser vytváří abstraktní syntaktický strom z uzlů \texttt{ASTNode}. Uzel má typ (\texttt{ASTNodeType}), odkaz \texttt{next} na dalšího sourozence a datovou část podle typu uzlu. Seznamy (příkazy bloku, parametry funkce, argumenty volání, funkce programu) jsou řetězeny přes \texttt{next}. Kořenem je \texttt{AST\_PROGRAM} se seznamem funkcí (\texttt{AST\_FUNC\_DEF}), tělo funkce, větve \texttt{if} a tělo \texttt{while} jsou bloky (\texttt{AST\_BLOCK}). Každý uzel nese datový typ hodnoty, kterou představuje (\texttt{expr\_type}); u definice proměnné a přiřazení je to typ proměnné, u funkce její návratový typ a u parametru jeho typ. Uzel funkce nese navíc příznak \texttt{is\_main}. Jednořádkové tělo funkce se reprezentuje jako blok s jediným výrazem.

\subsection{Tabulka symbolů}

Tabulku symbolů (\file{symtable.c}) jsme podle zvolené varianty implementovali jako vyvážený binární vyhledávací strom (AVL), jehož klíčem je jméno identifikátoru. Parser s ní pracuje jako se zásobníkem rozsahů: každý blok kódu otevírá nový rozsah (nový AVL strom) a při jeho opuštění se rozsah zruší. Rozsah s indexem~0 je globální a obsahuje funkce, včetně vestavěných, které se vkládají při inicializaci. Vyhledávání prochází rozsahy od nejvnitřnějšího, takže vnořená definice překryje stejnojmennou proměnnou z vnějšího bloku. Vložení jména do rozsahu, kde už je, je chyba redefinice (kód~4).

Záznam \texttt{symtable\_entry} obsahuje typ, příznaky (funkce, proměnná, nemodifikovatelná hodnota), pro funkce počet a typy parametrů a stav analýzy těla funkce (\texttt{BODY\_PENDING}, \texttt{BODY\_IN\_PROGRESS}, \texttt{BODY\_DONE}) s pozicemi v bufferu tokenů, které umožňují analyzovat tělo odloženě. Návratové kódy překladače jsou v tabulce~\ref{tab:errors}.

\begin{table}[ht]
\centering
\begin{tabular}{@{}cl@{}}
  \toprule
  Kód & Význam \\
  \midrule
  1 & lexikální chyba \\
  2 & syntaktická chyba (včetně chybějícího prologu) \\
  3 & nedefinovaná funkce, proměnná nebo typ; chybějící \ic{@main} \\
  4 & redefinice funkce, proměnné nebo typu \\
  5 & špatný počet nebo typ argumentů; nepovolené parametry \ic{@main} \\
  6 & typová nekompatibilita, nelze odvodit typ \\
  7 & zápis do nemodifikovatelné proměnné nebo parametru \\
  10 & ostatní sémantické chyby \\
  99 & interní chyba překladače \\
  \bottomrule
\end{tabular}
\caption{Návratové kódy překladače.}
\label{tab:errors}
\end{table}

\subsection{Generátor cílového kódu}

Generátor (\file{generator.c}) prochází AST a vypisuje kód IFJcode26. Veškerý jeho stav je v kontextu \texttt{GeneratorContext}, takže nepoužívá globální proměnné. Kód se skládá do bufferu v paměti a zapíše se na výstup až po úspěšném průchodu celým AST; při chybě nevznikne částečný výstup.

\subsubsection{Konvence volání funkcí a využití rámců}
\label{sec:konvence_volani}

Každá funkce má návěští \ic{\$fn\_}\textit{jméno}, začíná instrukcí \ic{PUSHFRAME} a končí \ic{POPFRAME} a \ic{RETURN}. Všechny proměnné funkce se definují najednou hned za \ic{PUSHFRAME}, jinak by se definice uvnitř těla cyklu provedla při každé iteraci znovu a interpret by skončil chybou. Počet definic zjistí generátor předem rekurzivním průchodem tělem funkce. Proměnná se jmenuje \ic{LF@\$v}\textit{n}, kde \textit{n} je pořadové číslo deklarace ve funkci (parametry mají čísla $0$ až $k-1$). Jména se přiřazují zásobníkem dvojic (jméno, číslo), který se na konci bloku zkrátí, čímž končí viditelnost proměnné; stejně pojmenované proměnné v různých blocích tak mají různá čísla.

\emph{Volající} nejprve vyhodnotí všechny argumenty na datový zásobník, pak vytvoří dočasný rámec (\ic{CREATEFRAME}), argumenty ze zásobníku odzadu uloží do proměnných \ic{TF@\$v}\textit{i} a zavolá funkci (\ic{CALL}). Vyhodnocení všech argumentů před vytvořením rámce je nutné kvůli rozšíření FUNEXP: vnořené volání by jinak přepsalo už rozpracovaný dočasný rámec. Volaná funkce předává návratovou hodnotu přes globální proměnnou \ic{GF@\_ret}, volající ji po návratu uloží na zásobník, kde se s ní dále pracuje jako s hodnotou výrazu. Hodnotou funkce je poslední výraz těla nebo hodnota příkazu \texttt{return}; ve funkci s návratovým typem \texttt{Unit} se hodnota zahodí a vrací se \ic{nil@nil}. Úvodní kód programu zavolá funkci s anotací \ic{@main} a skončí instrukcí \ic{EXIT int@0}. Globální rámec se kromě \ic{GF@\_ret} používá pro pomocné proměnné \ic{GF@\_t1} až \ic{GF@\_t6}.

\subsubsection{Vyhodnocování výrazů}

Výrazy se vyhodnocují zásobníkově: hodnota výrazu je vždy na vrcholu datového zásobníku a kód se generuje průchodem AST. Například sčítání proměnné \texttt{a} a konstanty \texttt{8} se přeloží na

\begin{center}
\texttt{PUSHS LF@\$v0\\PUSHS int@8\\ADDS}
\end{center}

Je-li jeden operand Int a druhý Double, převede se Int hned po vyhodnocení instrukcí \ic{INT2FLOATS}. Relace \texttt{<} a \texttt{>} odpovídají \ic{LTS} a \ic{GTS}; \texttt{<=} a \texttt{>=} se skládají z \ic{LTS}/\ic{GTS}, \ic{EQS} a \ic{ORS}, protože negace opačné relace by pro hodnotu NaN dávala špatný výsledek. Konkatenace řetězců nemá zásobníkovou instrukci, proto se hodnoty vyzvednou do pomocných proměnných a použije se \ic{CONCAT}. Podmínka příkazů \texttt{if} a \texttt{while} se vyhodnotí na zásobník, porovná s \ic{bool@false} a instrukcí \ic{JUMPIFEQS} se skočí za příslušnou větev. Návěští (\ic{\$else\_}\textit{n}, \ic{\$endif\_}\textit{n}, \ic{\$while\_}\textit{n}, \ic{\$endwhile\_}\textit{n}) jsou unikátní díky čítači v kontextu generátoru.

\subsubsection{Vestavěné funkce}

Jednoduché vestavěné funkce (\texttt{print}, \texttt{toStr}, \texttt{toInt}, \texttt{strLength}, \texttt{chr}, \texttt{readLine}) se generují přímo v místě volání podle statického typu argumentu. Složitější případy jsou pomocné rutiny ve výstupním kódu, volané instrukcí \ic{CALL}; argumenty dostávají na datovém zásobníku a výsledek na něm zanechávají (tabulka~\ref{tab:helpers}).

\begin{table}[ht]
\centering
\begin{tabularx}{\linewidth}{@{}lX@{}}
  \toprule
  Rutina & Účel \\
  \midrule
  \ic{\$\$idiv} & celočíselné dělení Int zaokrouhlené k nule; výsledek \ic{IDIV} se opraví podle znaménka zbytku, takže rutina funguje při zaokrouhlení k nule i k $-\infty$ \\
  \ic{\$\$ord} & \texttt{ord}; mimo rozsah vrací $-1$ \\
  \ic{\$\$substr} & \texttt{subStr}; při špatných mezích \ic{EXIT int@40} \\
  \ic{\$\$readint}, \ic{\$\$readdouble} & čtení čísla; při neplatném vstupu nebo EOF vrací 0 \\
  \bottomrule
\end{tabularx}
\caption{Pomocné rutiny generátoru.}
\label{tab:helpers}
\end{table}

Zadání uvádí, že \ic{IDIV} zaokrouhluje k $-\infty$, interpret \file{ic26int} ale dělí k nule; proto je \ic{\$\$idiv} napsaná nezávisle na zaokrouhlení. Instrukce \ic{READ} ukládá při chybném vstupu a na konci souboru hodnotu \ic{nil}, kdežto IFJ26 požaduje nulu; to řeší rutiny \ic{\$\$readint} a \ic{\$\$readdouble}. Dělení nulou a \texttt{chr} mimo rozsah hlásí přímo interpret.

\subsection{Přehled modulů}

\begin{tabularx}{\linewidth}{@{}lX@{}}
  \toprule
  Modul & Popis \\
  \midrule
  \file{main.c} & Čte zdrojový kód, spouští parser a po jeho úspěchu generátor, vrací návratový kód. \\
  \file{lexer.c/h} & Lexikální analyzátor: převádí vstup na tokeny, zpracovává odsazení a komentáře. \\
  \file{parser.c/h} & Syntaktický a sémantický analyzátor (rekurzivní sestup), staví AST. \\
  \file{precedence.c/h} & Precedenční analýza výrazů, typová kontrola výrazů a volání funkcí. \\
  \file{symtable.c/h} & Tabulka symbolů: zásobník rozsahů s AVL stromy a vestavěné funkce. \\
  \file{ast.c/h} & Uzly abstraktního syntaktického stromu a jejich tvorba. \\
  \file{generator.c/h} & Generátor kódu IFJcode26 z AST. \\
  \file{DynBuff.c/h} & Dynamický řetězcový buffer (lexer, generátor). \\
  \file{Error.h} & Výčet návratových kódů překladače. \\
  \bottomrule
\end{tabularx}

% =============================================================================
\section{Rozšíření}
% =============================================================================
Rozšíření jsou zpracována hlavně v parseru a generátor oproti základnímu zadání mění jen málo. \todo{u každého rozšíření upravit popis podle skutečné implementace}

\subsection{FUNEXP}
Volání vestavěné i uživatelské funkce je součástí výrazu a argumentem může být libovolný výraz. Parser volání zpracovává rekurzivně (\texttt{parse\_call}) a ve výrazu se chová jako operand. Generátor vyhodnocuje argumenty na datový zásobník před vytvořením rámce (viz \ref{sec:konvence_volani}), proto je vnořená volání nepřepisují.

\subsection{BOOLTHEN}
Typ \texttt{Boolean}, literály \texttt{true} a \texttt{false} a operátory \texttt{!}, \texttt{\&\&}, \texttt{||} jsou v precedenční analýze s prioritami podle zadání. Operátory \texttt{\&\&} a \texttt{||} se vyhodnocují zkratově instrukcí \ic{JUMPIFEQS}: pravý operand se nevyhodnotí, je-li výsledek znám z levého. Pro \texttt{toStr} nad hodnotou typu Boolean slouží pomocná rutina a funkce \texttt{readBoolean} převádí \ic{nil} (neplatný vstup) na \texttt{false}. LL gramatika dovoluje \texttt{if} bez \texttt{else} i řetězení \texttt{else if}.

\subsection{CONTROL}
Příkaz \texttt{if} lze použít jako výraz: každá větev nechá na zásobníku jednu hodnotu, převedenou na společný typ. Cyklus \texttt{for} s generátorem \texttt{to}, \texttt{until} a \texttt{by} a se stráží vyhodnotí meze a krok jednou do pomocných proměnných; nulový krok ukončí program s kódem~40. Pomocné proměnné cyklu se definují stejně jako ostatní na začátku funkce.

\subsection{VALMEM}
Zápis \texttt{výraz.člen} (\texttt{toString}, \texttt{toInt}, \texttt{floor}, \texttt{length}) parser převádí na volání vestavěných funkcí. Pro \texttt{floor} existuje pomocná rutina, protože IFJcode26 nemá instrukci pro zaokrouhlení dolů.

\subsection{INTERPOLATEDS}
Interpolovaný řetězec se rozloží na literály a vložené výrazy; každý výraz se převede na řetězec stejně jako ve funkci \texttt{toStr} a všechny části se spojí konkatenací.

\subsection{EXTKW}
Rozšiřující metoda se přeloží jako obyčejná funkce, jejímž prvním parametrem je příjemce. Parser jí podle typu příjemce a jména metody přidělí jedinečné jméno (např.\ \texttt{Int\$plus}), takže metody stejného jména pro různé typy nekolidují. Volání \texttt{x.plus(12)} se změní na volání takové funkce.

\subsection{PARAMLISTS}
Funkce s více seznamy parametrů se zploští na funkci s jedním spojeným seznamem parametrů a volání s více seznamy argumentů na jedno volání, argumenty se vyhodnocují zleva doprava. Funkce bez seznamu parametrů se volá jako funkce bez argumentů.

% =============================================================================
\section{Využití prostředků AI}
% =============================================================================
\todo{popsat podle požadavků sekce 4 hlavního zadání: jaké nástroje AI, k čemu (např.\ generátor kódu, testy, dokumentace), co z výstupu bylo převzato nebo upraveno a jak jsme výsledek ověřili (testy, ruční kontrola, srovnání se Scalou)}

% =============================================================================
\section*{Reference}
% =============================================================================
\begin{thebibliography}{9}
  \bibitem{zadani} Ondryáš, O., Křivka, Z.: \emph{Specifikace jazyků IFJ26 a IFJcode26 (IFJ/IAL 2026/2027)}. FIT VUT v Brně.
  \bibitem{scala} \emph{Scala 3 documentation}. \url{https://docs.scala-lang.org/scala3/}
  \bibitem{prednasky} Přednášky předmětů IFJ a IAL, FIT VUT v Brně (lexikální a syntaktická analýza, precedenční analýza, AVL stromy). \todo{přesné zdroje}
\end{thebibliography}
\todo{doplnit všechny převzaté zdroje včetně kódu a algoritmů, u převzatého kódu uvést místo použití}

% =============================================================================
\section*{Přílohy}
% =============================================================================
\subsection*{A\quad Diagram konečného automatu pro lexikální analýzu}
\label{sec:fa}
\begin{figure}[ht]
\centering
\fbox{\parbox{0.92\linewidth}{\centering\vspace{4cm}
  \todo{vložit diagram konečného automatu lexeru: stavy podle \texttt{enum State} v \file{lexer.h}, koncové stavy, přechody}
  \vspace{4cm}}}
\end{figure}
\newpage

\subsection*{B\quad LL gramatika pro rekurzivní sestup}
\label{sec:ll-gram}
Terminály \textit{INDENT} a \textit{DEDENT} jsou virtuální: lexer vrací počet mezer na začátku řádku, parser z něj odvozuje začátek a konec bloku. Neterminál \nt{expr} zpracovává precedenční analýza. Neterminály \nt{def\_head}, \nt{ext\_head} a \nt{else\_part} se jmenují jinak než stejnojmenná klíčová slova, aby se nepletly.

% Vygenerováno skriptem gen_tabulky.py, needitujte ručně.
{\renewcommand{\arraystretch}{1.12}\small
\begin{longtable}{@{}r@{\hspace{8pt}}r@{\hspace{5pt}}c@{\hspace{5pt}}l@{}}
1. & \nt{program} & $\longrightarrow$ & \nt{eols} \texttt{import} \texttt{ifj26} \texttt{.} \texttt{*} \nt{fun\_list} \\
2. & \nt{eols} & $\longrightarrow$ & \textit{EOL} \nt{eols} \\
3. & \nt{eols} & $\longrightarrow$ & $\varepsilon$ \\
4. & \nt{fun\_list} & $\longrightarrow$ & \nt{fun\_def} \nt{fun\_list} \\
5. & \nt{fun\_list} & $\longrightarrow$ & \nt{ext\_head} \nt{fun\_list} \\
6. & \nt{fun\_list} & $\longrightarrow$ & \textit{EOL} \nt{fun\_list} \\
7. & \nt{fun\_list} & $\longrightarrow$ & $\varepsilon$ \\
8. & \nt{fun\_def} & $\longrightarrow$ & \texttt{@main} \nt{eols} \nt{def\_head} \\
9. & \nt{fun\_def} & $\longrightarrow$ & \nt{def\_head} \\
10. & \nt{def\_head} & $\longrightarrow$ & \texttt{def} \textit{id} \nt{param\_lists} \nt{ret\_type} \texttt{=} \nt{fun\_body} \\
11. & \nt{param\_lists} & $\longrightarrow$ & \texttt{(} \nt{params} \texttt{)} \nt{param\_lists} \\
12. & \nt{param\_lists} & $\longrightarrow$ & $\varepsilon$ \\
13. & \nt{params} & $\longrightarrow$ & \nt{param} \nt{params\_n} \\
14. & \nt{params} & $\longrightarrow$ & $\varepsilon$ \\
15. & \nt{params\_n} & $\longrightarrow$ & \texttt{,} \nt{param} \nt{params\_n} \\
16. & \nt{params\_n} & $\longrightarrow$ & $\varepsilon$ \\
17. & \nt{param} & $\longrightarrow$ & \textit{id} \texttt{:} \nt{type} \\
18. & \nt{ret\_type} & $\longrightarrow$ & \texttt{:} \nt{type} \\
19. & \nt{ret\_type} & $\longrightarrow$ & $\varepsilon$ \\
20. & \nt{type} & $\longrightarrow$ & \texttt{Int} \\
21. & \nt{type} & $\longrightarrow$ & \texttt{Double} \\
22. & \nt{type} & $\longrightarrow$ & \texttt{String} \\
23. & \nt{type} & $\longrightarrow$ & \texttt{Unit} \\
24. & \nt{type} & $\longrightarrow$ & \texttt{Boolean} \\
25. & \nt{fun\_body} & $\longrightarrow$ & \nt{expr} \\
26. & \nt{fun\_body} & $\longrightarrow$ & \textit{EOL} \nt{block} \\
27. & \nt{ext\_head} & $\longrightarrow$ & \texttt{extension} \texttt{(} \textit{id} \texttt{:} \nt{type} \texttt{)} \textit{EOL} \textit{INDENT} \nt{ext\_defs} \textit{DEDENT} \\
28. & \nt{ext\_defs} & $\longrightarrow$ & \textit{EOL} \nt{ext\_defs} \\
29. & \nt{ext\_defs} & $\longrightarrow$ & \nt{def\_head} \nt{ext\_defs} \\
30. & \nt{ext\_defs} & $\longrightarrow$ & $\varepsilon$ \\
31. & \nt{block} & $\longrightarrow$ & \textit{EOL} \textit{INDENT} \nt{stmt} \nt{stmt\_list} \textit{DEDENT} \\
32. & \nt{stmt\_list} & $\longrightarrow$ & \textit{EOL} \nt{stmt\_rest} \\
33. & \nt{stmt\_list} & $\longrightarrow$ & \texttt{;} \nt{stmt\_rest} \\
34. & \nt{stmt\_list} & $\longrightarrow$ & $\varepsilon$ \\
35. & \nt{stmt\_rest} & $\longrightarrow$ & \nt{stmt} \nt{stmt\_list} \\
36. & \nt{stmt\_rest} & $\longrightarrow$ & \textit{EOL} \nt{stmt\_rest} \\
37. & \nt{stmt\_rest} & $\longrightarrow$ & $\varepsilon$ \\
38. & \nt{stmt} & $\longrightarrow$ & \texttt{val} \textit{id} \nt{opt\_type} \texttt{=} \nt{expr} \\
39. & \nt{stmt} & $\longrightarrow$ & \texttt{var} \textit{id} \nt{opt\_type} \texttt{=} \nt{expr} \\
40. & \nt{stmt} & $\longrightarrow$ & \textit{id} \texttt{=} \nt{expr} \\
41. & \nt{stmt} & $\longrightarrow$ & \nt{expr} \\
42. & \nt{stmt} & $\longrightarrow$ & \texttt{return} \nt{ret\_val} \\
43. & \nt{stmt} & $\longrightarrow$ & \texttt{if} \nt{expr} \texttt{then} \nt{body} \nt{else\_part} \\
44. & \nt{stmt} & $\longrightarrow$ & \texttt{while} \nt{expr} \texttt{do} \nt{body} \\
45. & \nt{stmt} & $\longrightarrow$ & \texttt{for} \nt{for\_head} \texttt{do} \nt{body} \\
46. & \nt{opt\_type} & $\longrightarrow$ & \texttt{:} \nt{type} \\
47. & \nt{opt\_type} & $\longrightarrow$ & $\varepsilon$ \\
48. & \nt{ret\_val} & $\longrightarrow$ & \nt{expr} \\
49. & \nt{ret\_val} & $\longrightarrow$ & $\varepsilon$ \\
50. & \nt{else\_part} & $\longrightarrow$ & \texttt{else} \nt{body} \\
51. & \nt{else\_part} & $\longrightarrow$ & \textit{EOL} \texttt{else} \nt{body} \\
52. & \nt{else\_part} & $\longrightarrow$ & $\varepsilon$ \\
53. & \nt{body} & $\longrightarrow$ & \nt{block} \\
54. & \nt{body} & $\longrightarrow$ & \nt{stmt} \\
55. & \nt{for\_head} & $\longrightarrow$ & \nt{generator} \nt{step} \nt{guards} \\
56. & \nt{for\_head} & $\longrightarrow$ & \textit{EOL} \textit{INDENT} \nt{generator} \nt{step} \nt{guards\_ml} \textit{DEDENT} \\
57. & \nt{generator} & $\longrightarrow$ & \textit{id} \texttt{\textless-} \nt{expr} \nt{range} \nt{expr} \\
58. & \nt{range} & $\longrightarrow$ & \texttt{to} \\
59. & \nt{range} & $\longrightarrow$ & \texttt{until} \\
60. & \nt{step} & $\longrightarrow$ & \texttt{by} \nt{expr} \\
61. & \nt{step} & $\longrightarrow$ & $\varepsilon$ \\
62. & \nt{guards} & $\longrightarrow$ & \texttt{if} \nt{expr} \nt{guards} \\
63. & \nt{guards} & $\longrightarrow$ & $\varepsilon$ \\
64. & \nt{guards\_ml} & $\longrightarrow$ & \textit{EOL} \texttt{if} \nt{expr} \nt{guards\_ml} \\
65. & \nt{guards\_ml} & $\longrightarrow$ & $\varepsilon$ \\
66. & \nt{if\_expr} & $\longrightarrow$ & \texttt{if} \nt{expr} \texttt{then} \nt{branch} \nt{else\_e} \\
67. & \nt{else\_e} & $\longrightarrow$ & \texttt{else} \nt{branch} \\
68. & \nt{else\_e} & $\longrightarrow$ & \textit{EOL} \texttt{else} \nt{branch} \\
69. & \nt{branch} & $\longrightarrow$ & \nt{expr} \\
70. & \nt{branch} & $\longrightarrow$ & \nt{block} \\
\end{longtable}}

\clearpage

\begin{landscape}
\thispagestyle{empty}
\subsection*{C\quad LL tabulka}
\label{sec:ll-tab}
\begin{center}
\resizebox{0.8\linewidth}{!}{% Vygenerováno skriptem gen_tabulky.py, needitujte ručně.
{\scriptsize\setlength{\tabcolsep}{2.2pt}
\renewcommand{\arraystretch}{1.08}
\begin{tabular}{|l|c|c|c|c|c|c|c|c|c|c|c|c|c|c|c|c|c|c|c|c|c|c|c|c|c|c|c|c|c|c|c|c|c|c|c|}
\hline
& \rotatebox{90}{\texttt{import}\hspace{2pt}} & \rotatebox{90}{\texttt{@main}\hspace{2pt}} & \rotatebox{90}{\texttt{def}\hspace{2pt}} & \rotatebox{90}{\texttt{extension}\hspace{2pt}} & \rotatebox{90}{\textit{id}\hspace{2pt}} & \rotatebox{90}{\texttt{(}\hspace{2pt}} & \rotatebox{90}{\texttt{)}\hspace{2pt}} & \rotatebox{90}{\texttt{,}\hspace{2pt}} & \rotatebox{90}{\texttt{:}\hspace{2pt}} & \rotatebox{90}{\texttt{=}\hspace{2pt}} & \rotatebox{90}{\texttt{;}\hspace{2pt}} & \rotatebox{90}{\texttt{Int}\hspace{2pt}} & \rotatebox{90}{\texttt{Double}\hspace{2pt}} & \rotatebox{90}{\texttt{String}\hspace{2pt}} & \rotatebox{90}{\texttt{Unit}\hspace{2pt}} & \rotatebox{90}{\texttt{Boolean}\hspace{2pt}} & \rotatebox{90}{\textit{literal}\hspace{2pt}} & \rotatebox{90}{\texttt{true}\hspace{2pt}} & \rotatebox{90}{\texttt{false}\hspace{2pt}} & \rotatebox{90}{\texttt{-}\hspace{2pt}} & \rotatebox{90}{\texttt{!}\hspace{2pt}} & \rotatebox{90}{\texttt{val}\hspace{2pt}} & \rotatebox{90}{\texttt{var}\hspace{2pt}} & \rotatebox{90}{\texttt{return}\hspace{2pt}} & \rotatebox{90}{\texttt{if}\hspace{2pt}} & \rotatebox{90}{\texttt{else}\hspace{2pt}} & \rotatebox{90}{\texttt{while}\hspace{2pt}} & \rotatebox{90}{\texttt{do}\hspace{2pt}} & \rotatebox{90}{\texttt{for}\hspace{2pt}} & \rotatebox{90}{\texttt{to}\hspace{2pt}} & \rotatebox{90}{\texttt{until}\hspace{2pt}} & \rotatebox{90}{\texttt{by}\hspace{2pt}} & \rotatebox{90}{\textit{DEDENT}\hspace{2pt}} & \rotatebox{90}{\textit{EOL}\hspace{2pt}} & \rotatebox{90}{\textit{EOF}\hspace{2pt}} \\ \hline
\nt{program} & 1 &  &  &  &  &  &  &  &  &  &  &  &  &  &  &  &  &  &  &  &  &  &  &  &  &  &  &  &  &  &  &  &  & 1 &  \\ \hline
\nt{eols} & 3 &  & 3 &  &  &  &  &  &  &  &  &  &  &  &  &  &  &  &  &  &  &  &  &  &  &  &  &  &  &  &  &  &  & 2 &  \\ \hline
\nt{fun\_list} &  & 4 & 4 & 5 &  &  &  &  &  &  &  &  &  &  &  &  &  &  &  &  &  &  &  &  &  &  &  &  &  &  &  &  &  & 6 & 7 \\ \hline
\nt{fun\_def} &  & 8 & 9 &  &  &  &  &  &  &  &  &  &  &  &  &  &  &  &  &  &  &  &  &  &  &  &  &  &  &  &  &  &  &  &  \\ \hline
\nt{def\_head} &  &  & 10 &  &  &  &  &  &  &  &  &  &  &  &  &  &  &  &  &  &  &  &  &  &  &  &  &  &  &  &  &  &  &  &  \\ \hline
\nt{param\_lists} &  &  &  &  &  & 11 &  &  & 12 & 12 &  &  &  &  &  &  &  &  &  &  &  &  &  &  &  &  &  &  &  &  &  &  &  &  &  \\ \hline
\nt{params} &  &  &  &  & 13 &  & 14 &  &  &  &  &  &  &  &  &  &  &  &  &  &  &  &  &  &  &  &  &  &  &  &  &  &  &  &  \\ \hline
\nt{params\_n} &  &  &  &  &  &  & 16 & 15 &  &  &  &  &  &  &  &  &  &  &  &  &  &  &  &  &  &  &  &  &  &  &  &  &  &  &  \\ \hline
\nt{param} &  &  &  &  & 17 &  &  &  &  &  &  &  &  &  &  &  &  &  &  &  &  &  &  &  &  &  &  &  &  &  &  &  &  &  &  \\ \hline
\nt{ret\_type} &  &  &  &  &  &  &  &  & 18 & 19 &  &  &  &  &  &  &  &  &  &  &  &  &  &  &  &  &  &  &  &  &  &  &  &  &  \\ \hline
\nt{type} &  &  &  &  &  &  &  &  &  &  &  & 20 & 21 & 22 & 23 & 24 &  &  &  &  &  &  &  &  &  &  &  &  &  &  &  &  &  &  &  \\ \hline
\nt{fun\_body} &  &  &  &  & 25 & 25 &  &  &  &  &  &  &  &  &  &  & 25 & 25 & 25 & 25 & 25 &  &  &  & 25 &  &  &  &  &  &  &  &  & 26 &  \\ \hline
\nt{ext\_head} &  &  &  & 27 &  &  &  &  &  &  &  &  &  &  &  &  &  &  &  &  &  &  &  &  &  &  &  &  &  &  &  &  &  &  &  \\ \hline
\nt{ext\_defs} &  &  & 29 &  &  &  &  &  &  &  &  &  &  &  &  &  &  &  &  &  &  &  &  &  &  &  &  &  &  &  &  &  & 30 & 28 &  \\ \hline
\nt{block} &  &  &  &  &  &  &  &  &  &  &  &  &  &  &  &  &  &  &  &  &  &  &  &  &  &  &  &  &  &  &  &  &  & 31 &  \\ \hline
\nt{stmt\_list} &  &  &  &  &  &  &  &  &  &  & 33 &  &  &  &  &  &  &  &  &  &  &  &  &  &  &  &  &  &  &  &  &  & 34 & 32 &  \\ \hline
\nt{stmt\_rest} &  &  &  &  & 35 & 35 &  &  &  &  &  &  &  &  &  &  & 35 & 35 & 35 & 35 & 35 & 35 & 35 & 35 & 35 &  & 35 &  & 35 &  &  &  & 37 & 36 &  \\ \hline
\nt{stmt} &  &  &  &  & \cf{40/41} & 41 &  &  &  &  &  &  &  &  &  &  & 41 & 41 & 41 & 41 & 41 & 38 & 39 & 42 & \cf{41/43} &  & 44 &  & 45 &  &  &  &  &  &  \\ \hline
\nt{opt\_type} &  &  &  &  &  &  &  &  & 46 & 47 &  &  &  &  &  &  &  &  &  &  &  &  &  &  &  &  &  &  &  &  &  &  &  &  &  \\ \hline
\nt{ret\_val} &  &  &  &  & 48 & 48 &  &  &  &  & 49 &  &  &  &  &  & 48 & 48 & 48 & 48 & 48 &  &  &  & 48 & 49 &  &  &  &  &  &  & 49 & 49 &  \\ \hline
\nt{else\_part} &  &  &  &  &  &  &  &  &  &  & 52 &  &  &  &  &  &  &  &  &  &  &  &  &  &  & \cf{50/52} &  &  &  &  &  &  & 52 & \cf{51/52} &  \\ \hline
\nt{body} &  &  &  &  & 54 & 54 &  &  &  &  &  &  &  &  &  &  & 54 & 54 & 54 & 54 & 54 & 54 & 54 & 54 & 54 &  & 54 &  & 54 &  &  &  &  & 53 &  \\ \hline
\nt{for\_head} &  &  &  &  & 55 &  &  &  &  &  &  &  &  &  &  &  &  &  &  &  &  &  &  &  &  &  &  &  &  &  &  &  &  & 56 &  \\ \hline
\nt{generator} &  &  &  &  & 57 &  &  &  &  &  &  &  &  &  &  &  &  &  &  &  &  &  &  &  &  &  &  &  &  &  &  &  &  &  &  \\ \hline
\nt{range} &  &  &  &  &  &  &  &  &  &  &  &  &  &  &  &  &  &  &  &  &  &  &  &  &  &  &  &  &  & 58 & 59 &  &  &  &  \\ \hline
\nt{step} &  &  &  &  &  &  &  &  &  &  &  &  &  &  &  &  &  &  &  &  &  &  &  &  & 61 &  &  & 61 &  &  &  & 60 & 61 & 61 &  \\ \hline
\nt{guards} &  &  &  &  &  &  &  &  &  &  &  &  &  &  &  &  &  &  &  &  &  &  &  &  & 62 &  &  & 63 &  &  &  &  &  &  &  \\ \hline
\nt{guards\_ml} &  &  &  &  &  &  &  &  &  &  &  &  &  &  &  &  &  &  &  &  &  &  &  &  &  &  &  &  &  &  &  &  & 65 & 64 &  \\ \hline
\nt{if\_expr} &  &  &  &  &  &  &  &  &  &  &  &  &  &  &  &  &  &  &  &  &  &  &  &  & 66 &  &  &  &  &  &  &  &  &  &  \\ \hline
\nt{else\_e} &  &  &  &  &  &  &  &  &  &  &  &  &  &  &  &  &  &  &  &  &  &  &  &  &  & 67 &  &  &  &  &  &  &  & 68 &  \\ \hline
\nt{branch} &  &  &  &  & 69 & 69 &  &  &  &  &  &  &  &  &  &  & 69 & 69 & 69 & 69 & 69 &  &  &  & 69 &  &  &  &  &  &  &  &  & 70 &  \\ \hline
\end{tabular}}
}
\end{center}
\par\small
\textbf{Poznámky.} Sloupec \textit{literal} zahrnuje celočíselné, desetinné i řetězcové literály. Neterminál \nt{expr} zpracovává precedenční analýza, proto pro něj tabulka nemá řádek; FIRST(\nt{expr}) $=\{$\textit{id}, \textit{literal}, \texttt{true}, \texttt{false}, \texttt{(}, \texttt{-}, \texttt{!}, \texttt{if}$\}$. Tokeny \texttt{to}, \texttt{until} a \texttt{by} jsou identifikátory se zadaným lexémem a \texttt{<-} vrací lexer jako dvojici tokenů \texttt{<} a \texttt{-}. Sloupce terminálů, které nejsou v žádné buňce (\texttt{ifj26}, \texttt{.}, \texttt{*}, \texttt{<-}, \texttt{then}, \textit{INDENT}), jsou vynechány.

\textbf{Konflikty} (červeně) a jejich řešení mimo tabulku:
\begin{itemize}[itemsep=1pt]
  \item[] \cf{40/41} (\textit{id}): přiřazení, je-li druhým tokenem \texttt{=}, jinak výraz.
  \item[] \cf{41/43} (\texttt{if}): na začátku příkazu se \texttt{if} zpracuje jako příkaz (pravidlo 43); \texttt{if} jako výraz (pravidlo 66) se použije na místě operandu. Je-li \texttt{if} posledním příkazem bloku, jehož hodnota se využije (např.\ tělo funkce), musí ho parser chápat jako výraz. \todo{ověřit, jak to parser řeší}
  \item[] \cf{50/52} (\texttt{else}): \texttt{else} patří nejbližšímu \texttt{if}.
  \item[] \cf{51/52} (\textit{EOL}): pravidlo 51, začíná-li další řádek klíčovým slovem \texttt{else} zarovnaným s \texttt{if}, jinak pravidlo 52.
\end{itemize}
\end{landscape}

\subsection*{D\quad Gramatika pro precedenční analýzu}
\label{sec:psa-gram}
% Vygenerováno skriptem gen_tabulky.py, needitujte ručně.
{\renewcommand{\arraystretch}{1.1}\small
\begin{tabular}{@{}r@{\hspace{8pt}}r@{\hspace{5pt}}c@{\hspace{5pt}}l@{}}
1. & $E$ & $\longrightarrow$ & \texttt{!} $E$ \\
2. & $E$ & $\longrightarrow$ & \texttt{-} $E$ \\
3. & $E$ & $\longrightarrow$ & \texttt{+} $E$ \\
4. & $E$ & $\longrightarrow$ & $E$ \texttt{*} $E$ \\
5. & $E$ & $\longrightarrow$ & $E$ \texttt{/} $E$ \\
6. & $E$ & $\longrightarrow$ & $E$ \texttt{+} $E$ \\
7. & $E$ & $\longrightarrow$ & $E$ \texttt{-} $E$ \\
8. & $E$ & $\longrightarrow$ & $E$ \texttt{<} $E$ \\
9. & $E$ & $\longrightarrow$ & $E$ \texttt{<=} $E$ \\
10. & $E$ & $\longrightarrow$ & $E$ \texttt{>} $E$ \\
11. & $E$ & $\longrightarrow$ & $E$ \texttt{>=} $E$ \\
12. & $E$ & $\longrightarrow$ & $E$ \texttt{==} $E$ \\
13. & $E$ & $\longrightarrow$ & $E$ \texttt{!=} $E$ \\
14. & $E$ & $\longrightarrow$ & $E$ \texttt{\&\&} $E$ \\
15. & $E$ & $\longrightarrow$ & $E$ \texttt{\textbar\textbar} $E$ \\
16. & $E$ & $\longrightarrow$ & \texttt{(} $E$ \texttt{)} \\
17. & $E$ & $\longrightarrow$ & \textit{i} \\
18. & $E$ & $\longrightarrow$ & \textit{f} \texttt{(} \texttt{)} \\
19. & $E$ & $\longrightarrow$ & \textit{f} \texttt{(} $E$ \texttt{)} \\
20. & $E$ & $\longrightarrow$ & \textit{f} \texttt{(} $E$ \texttt{,} \ldots \texttt{)} \\
\end{tabular}}


\subsection*{E\quad Precedenční tabulka}
\label{sec:prec-tab}
Priority a asociativita operátorů podle zadání:

\begin{center}
% Vygenerováno skriptem gen_tabulky.py, needitujte ručně.
{\renewcommand{\arraystretch}{1.15}
\begin{tabular}{@{}cll@{}}
\toprule Priorita & Operátory & Asociativita \\ \midrule
1 & \texttt{+}$_{\mathrm{un}}$ \texttt{-}$_{\mathrm{un}}$ \texttt{!} & pravá \\
2 & \texttt{*}\ \texttt{/} & levá \\
3 & \texttt{+}\ \texttt{-} & levá \\
4 & \texttt{<}\ \texttt{<=}\ \texttt{>}\ \texttt{>=} & žádná \\
5 & \texttt{==}\ \texttt{!=} & žádná \\
6 & \texttt{\&\&} & levá \\
7 & \texttt{\textbar\textbar} & levá \\
\bottomrule\end{tabular}}

\end{center}

Symboly v tabulce: $\lessdot$ posun, $\gtrdot$ redukce, $\doteq$ shoda závorek, A přijetí, prázdná buňka chyba (včetně zakázaného řetězení relačních operátorů). Řádky odpovídají operátoru na vrcholu zásobníku, sloupce vstupnímu symbolu; \textit{i} je operand (literál, identifikátor, výsledek volání funkce).

\begin{center}
% Vygenerováno skriptem gen_tabulky.py, needitujte ručně.
{\small\setlength{\tabcolsep}{4pt}
\renewcommand{\arraystretch}{1.3}
\begin{tabular}{|l||c|c|c|c|c|c|c|c|c|c|c|}
\hline
 & \texttt{+}$_{\mathrm{un}}$ \texttt{-}$_{\mathrm{un}}$ \texttt{!} & \texttt{*}\ \texttt{/} & \texttt{+}\ \texttt{-} & \texttt{<}\ \texttt{<=}\ \texttt{>}\ \texttt{>=} & \texttt{==}\ \texttt{!=} & \texttt{\&\&} & \texttt{\textbar\textbar} & \texttt{(} & \texttt{)} & \textit{i} & \texttt{\$} \\ \hline\hline
\texttt{+}$_{\mathrm{un}}$ \texttt{-}$_{\mathrm{un}}$ \texttt{!} & $\lessdot$ & $\gtrdot$ & $\gtrdot$ & $\gtrdot$ & $\gtrdot$ & $\gtrdot$ & $\gtrdot$ & $\lessdot$ & $\gtrdot$ & $\lessdot$ & $\gtrdot$ \\ \hline
\texttt{*}\ \texttt{/} & $\lessdot$ & $\gtrdot$ & $\gtrdot$ & $\gtrdot$ & $\gtrdot$ & $\gtrdot$ & $\gtrdot$ & $\lessdot$ & $\gtrdot$ & $\lessdot$ & $\gtrdot$ \\ \hline
\texttt{+}\ \texttt{-} & $\lessdot$ & $\lessdot$ & $\gtrdot$ & $\gtrdot$ & $\gtrdot$ & $\gtrdot$ & $\gtrdot$ & $\lessdot$ & $\gtrdot$ & $\lessdot$ & $\gtrdot$ \\ \hline
\texttt{<}\ \texttt{<=}\ \texttt{>}\ \texttt{>=} & $\lessdot$ & $\lessdot$ & $\lessdot$ &  & $\gtrdot$ & $\gtrdot$ & $\gtrdot$ & $\lessdot$ & $\gtrdot$ & $\lessdot$ & $\gtrdot$ \\ \hline
\texttt{==}\ \texttt{!=} & $\lessdot$ & $\lessdot$ & $\lessdot$ & $\lessdot$ &  & $\gtrdot$ & $\gtrdot$ & $\lessdot$ & $\gtrdot$ & $\lessdot$ & $\gtrdot$ \\ \hline
\texttt{\&\&} & $\lessdot$ & $\lessdot$ & $\lessdot$ & $\lessdot$ & $\lessdot$ & $\gtrdot$ & $\gtrdot$ & $\lessdot$ & $\gtrdot$ & $\lessdot$ & $\gtrdot$ \\ \hline
\texttt{\textbar\textbar} & $\lessdot$ & $\lessdot$ & $\lessdot$ & $\lessdot$ & $\lessdot$ & $\lessdot$ & $\gtrdot$ & $\lessdot$ & $\gtrdot$ & $\lessdot$ & $\gtrdot$ \\ \hline
\texttt{(} & $\lessdot$ & $\lessdot$ & $\lessdot$ & $\lessdot$ & $\lessdot$ & $\lessdot$ & $\lessdot$ & $\lessdot$ & $\doteq$ & $\lessdot$ &  \\ \hline
\texttt{)} &  & $\gtrdot$ & $\gtrdot$ & $\gtrdot$ & $\gtrdot$ & $\gtrdot$ & $\gtrdot$ &  & $\gtrdot$ &  & $\gtrdot$ \\ \hline
\textit{i} &  & $\gtrdot$ & $\gtrdot$ & $\gtrdot$ & $\gtrdot$ & $\gtrdot$ & $\gtrdot$ &  & $\gtrdot$ &  & $\gtrdot$ \\ \hline
\texttt{\$} & $\lessdot$ & $\lessdot$ & $\lessdot$ & $\lessdot$ & $\lessdot$ & $\lessdot$ & $\lessdot$ & $\lessdot$ &  & $\lessdot$ & A \\ \hline
\end{tabular}}

\end{center}

\end{document}
